
\begin{frame}{Disaggregated Evaluation: The Core Technique}

    \textit{One number for the test set, then the same number for every slice you can define.}

    \vspace{0.7em}
    \centering
    \scriptsize
    \begin{tabular}{lrrrrr}
        \toprule
        \textbf{Slice} & \textbf{$n$} & \textbf{Selection rate} & \textbf{TPR} & \textbf{FPR} & \textbf{PPV} \\
        \midrule
        Overall            & 10{,}000 & 0.24 & 0.78 & 0.11 & 0.71 \\
        \midrule
        Group A            &  8{,}200 & 0.26 & 0.81 & 0.10 & 0.73 \\
        Group B            &  1{,}800 & 0.14 & 0.62 & 0.13 & 0.58 \\
        \midrule
        Group B, age $<30$ &    420   & 0.11 & 0.49 & 0.14 & 0.51 \\
        Group B, age $\geq 60$ &  180 & 0.09 & 0.44 & 0.12 & 0.47 \\
        \bottomrule
    \end{tabular}

    \vspace{0.5em}
    {\scriptsize Illustrative numbers, constructed for this slide.}

    \vspace{0.8em}
    \raggedright
    \begin{itemize}\scriptsize
        \setlength\itemsep{0.35em}
        \item The overall row is fine. Group B is not, and the intersectional rows are worse ---
              the pattern from Gender Shades, reproduced in miniature.
        \item Always print \textbf{$n$} beside the metric. Always print the \textbf{worst slice}
              in the summary, not just the mean.
        \item In \texttt{scikit-learn}, this is a \texttt{groupby} plus
              \texttt{classification\_report}; in \texttt{fairlearn} it is \texttt{MetricFrame}.
    \end{itemize}
\end{frame}
